\chapter*{Abstract}
\addcontentsline{toc}{chapter}{Abstract}
Nuclear fragmentation processes induced by charged particle beams represent a critical scientific frontier for both particle therapy in oncology and radiation shielding design in human space exploration. The FragmentatiOn Of Target (FOOT) experiment, approved by the Istituto Nazionale di Fisica Nucleare (INFN), was designed to measure double-differential nuclear fragmentation cross sections with unprecedented accuracy—targeting uncertainties below 10\% for target fragmentation and below 5\% for projectile fragmentation. To overcome the experimental challenge of detecting short-range (10–100 $\mu\text{m}$) target fragments, FOOT employs an inverse kinematics approach. This thesis focuses on the hardware, firmware, and software optimizations of the FOOT tracking system, aiming to maximize acquisition rates, refine detector alignment, and optimize spatial and momentum reconstruction.

On the hardware and firmware front, this work contributes to the upgrade of the Vertex Tracker (VT) detector, transitioning from MIMOSA-28 sensors to advanced MIMOSIS-2.1 Monolithic Active Pixel Sensors (MAPS). MIMOSIS-2.1 features a global shutter readout that drastically reduces frame readout time from 185.6 $\mu\text{s}$ down to 5 $\mu\text{s}$. To handle the associated high-rate Trigger and Data Acquisition (TDAQ) demands, a dedicated Trigger Board (TrgBox) firmware was developed in VHDL and implemented on an Intel Cyclone V System-on-Chip (SoC) FPGA hosted on a Terasic DE10-Nano board. The firmware incorporates a First-Word Fall-Through (FWFT) configurable FIFO readout logic, enhanced busy signal management, and C++ driver software interfacing the FPGA registers with the central DAQ. Validation tests performed at the Beam Test Facility (BTF) in Frascati demonstrated stable clock distribution, zero event loss, and reliable operation up to trigger rates of 20 kHz.

On the software side, offline track reconstruction algorithms within the SHOE (Software for Hadrontherapy Optimization Experiment) framework were systematically optimized. Using experimental data collected during the CNAO 2025 campaign with a 200 MeV/u $^{12}\text{C}$ beam impinging on a 5 mm graphite target, detailed pull distribution analyses were conducted for the MIMOSA-28 VT and the Microstrip Silicon Detector (MSD). A dedicated $\chi^2$ minimization algorithm was developed to identify and correct physical longitudinal ($z$-axis) misalignments in the VT planes, leading to an updated geometry model in SHOE. Furthermore, to resolve position error overestimation in the analog readout of the MSD, a non-linear position-finding algorithm based on Turchetta’s integral $f(\eta)$ method was implemented, utilizing a sum of error functions (erf) to linearize charge sharing across the 150 $\mu\text{m}$ readout pitch.

The combined hardware and software enhancements yielded a marked improvement in global track fitting quality. The implementation of the $f(\eta)$ algorithm suppressed unphysical position uncertainties, shifted the reduced $\chi^2$ peak of global track fits from over-estimated values to 0.86, and significantly flattened the $p$-value distribution by removing poorly fitted track artifacts. Ultimately, evaluating the reconstructed momentum using GENFIT and a 10-iteration forward-backward Kalman Filter demonstrated momentum resolutions of approximately 3.34\% for $^{12}\text{C}$ primary ions and 2.96\% for $^4\text{He}$ secondary fragments. These results fully satisfy the 3–5\% momentum precision requirement of the FOOT experiment, laying the foundation for high-precision cross section measurements in upcoming data campaigns.
